
%%%%%%%%%%%%%%%%%%%%%%%%%%%%%%%%%%%%%%%%
%           main body
%%%%%%%%%%%%%%%%%%%%%%%%%%%%%%%%%%%%%%%%


%-----------------------------------
%	TITLE PAGE
%-----------------------------------

\title[Probability \& Statistics]{\LARGE \bfseries 第三部分\quad 概率与统计} % The short title appears at the bottom of every slide, the full title is only on the title page
\author[Basic Concepts] % Your institution as it will appear on the bottom of every slide, may be shorthand to save space
{\Large \bfseries \S 1. 概率与统计基本概念}
% \author{Singular Value Decomposition} % Your name
% \date{\today} % Date, can be changed to a custom date
\date{}

%------------------------------------------------

\begin{document}


%------------------------------------------------

\begin{frame}

\titlepage % Print the title page as the first slide

\end{frame}

%------------------------------------------------

\begin{frame}

\begin{itemize}
\item 机器学习的问题很多时候就是一个概率统计的问题。
\vspace{1em}
\item 很多实际问题是带有不确定性的，需要用统计的方法找出规律，用概率的语言加以描述。
\end{itemize}

\end{frame}

%------------------------------------------------

\begin{frame}
\frametitle{内容提要} % Table of contents slide, comment this block out to remove it
\tableofcontents % Throughout your presentation, if you choose to use \section{} and \subsection{} commands, these will automatically be printed on this slide as an overview of your presentation
\end{frame}

%-----------------------------------
%	PRESENTATION SLIDES
%-----------------------------------

\section{随机事件}

%------------------------------------------------

\begin{frame}

\begin{block}{随机试验，Random Experiment}
随机试验指的是具有以下三个特点的试验：
\begin{itemize}
\item 可重复性：相同条件下可重复。
\item 可观测性：每次试验结果不止一个，并且能事先明确试验的所有可能的结果。
\item 不确定性：一次试验之前，不能预知会出现哪个结果。
\end{itemize}
\end{block}

\begin{block}{样本空间，Sample Space}
每次试验的每一个结果（Outcome）被称为基本事件（Elementary Event），或者样本点（Sample Point），通常被记作$\omega$, 全部样本点构成的集合，被称作样本空间，一般记作$\Omega$.
\end{block}

\end{frame}

%------------------------------------------------

\begin{frame}

\begin{block}{离散样本空间与连续样本空间}
\begin{itemize}
\item 一枚硬币被抛两次的样本空间
$$\Omega = \{\text{正正，正反，反正，反反}\}$$
\item 一个骰子投一次的点数的样本空间
$$\Omega = \{1,2,3,4,5,6\}$$
\pause
\item 某地气温的观测值的样本空间
$$\Omega = \{T \ |\ T > 0K\}$$
\item 某地房屋价格的样本空间
$$\Omega = \{p \ |\ p > 0\text{元}\}$$
\end{itemize}
\end{block}

\end{frame}

%------------------------------------------------

\begin{frame}

\begin{block}{随机事件，Random Event}
样本空间中，具有某些特征的样本点的集合叫做一个随机事件，常用$A,B,\ldots$表示。随机事件也简称事件。一个随机事件就是样本空间的一个子集。

一次试验中，如果出现了事件$A$包含的样本点，则称事件$A$发生。

\vspace{1em}
\pause

\begin{itemize}
\item 必然事件：必然发生的随机事件$A = \Omega$
\item 不可能事件：不可能发生的事件$A = \emptyset$
\end{itemize}
必然事件与不可能事件本质上不具有``不确定性''，但是为了方便讨论，依然将其视作``特殊''的随机事件。

\vspace{1em}
\pause

某地气温$\Omega = \{T \ |\ T > 0K\}$, 那么
\begin{itemize}
\item 气温高于$0K$就是一个必然事件；
\item 气温低于$0K$是不可能事件；
\item 气温高于$273K$（即约0摄氏度）是某随机事件。
\end{itemize}
\end{block}

\end{frame}

%------------------------------------------------

\begin{frame}

\begin{block}{随机事件的关系与运算}
\begin{itemize}
\item 事件包含：$A\subset B$
\item 事件的和（并）：$A+B = A\cup B$
\item 事件的积（交）：$AB = A\cap B$
\item 事件的差：$A-B$
\item 事件互斥：$A\cap B = \emptyset,$ 事件$A,B$不能同时发生
\item 事件对立：$A\cap B = \emptyset$且$A\cup B = \Omega$, 事件$A,B$不能同时发生，且恰好发生一个。
\end{itemize}
\end{block}

\end{frame}

%------------------------------------------------

\section{概率空间}

%------------------------------------------------

\subsection{古典概型}

%------------------------------------------------

\begin{frame}

\begin{block}{古典概型}
若基本事件都等可能出现，那么此时某随机事件$A$的概率为
\begin{align*}
& P(A) = \dfrac{\text{事件$A$包含基本事件数目}}{\text{样本空间包含基本事件数目}} \\
\text{或\hspace{1em}} & P(A) = \dfrac{\text{事件$A$所占区间大小}}{\text{样本空间所占区间大小}}
\end{align*}
\end{block}

\pause

\begin{block}{蒲丰投针问题，Buffon's Needle Problem}
设有一个以平行且等距木纹铺成的地板，现在随意抛一支长度$\ell$比木纹之间距离$h$小的针，那么针与纹线相交的概率为
$$\dfrac{2\ell}{\pi h}$$
\end{block}

\end{frame}

%------------------------------------------------

\begin{frame}

\begin{block}{伯努利（Bernoulli）概型}
如果随机试验只有两种可能的结果：事件$A$发生或不发生，则称这样的试验为伯努利试验。事件$A$发生的概率$p$被称作该伯努利试验的参数。

\vspace{1em}

将伯努利试验在相同条件下{\color{red}独立重复}地进行$n$次，这样的试验被称作$n$重伯努利试验，或$n$重独立重复试验，也简称为伯努利概型。
\end{block}

\end{frame}

%------------------------------------------------

\subsection{公理化概率}

%------------------------------------------------

\begin{frame}

\begin{block}{公理化概率}
设样本空间为$\Omega$, $\mathcal{F}$为$\Omega$某些子集（随机事件）组成的集合，满足
\begin{itemize}
\item $\Omega, \emptyset \in \mathcal{F}$
\item $A_1,A_2\in \mathcal{F} \Longrightarrow A_1\cup A_2 \in \mathcal{F}$
\item $A\in \mathcal{F} \Longrightarrow \Omega - A \in \mathcal{F}$
\end{itemize}
称作事件域。如果定义在$\mathcal{F}$上的一个实值函数$P: \mathcal{F} \to \mathbb{R}$满足
\begin{itemize}
\item 非负性：对任一随机事件$A\in\mathcal{F}$, $0 \le P(A) \le 1$
\item 正则性：$P(\Omega) = 1$
\item 可数可加性：若可数个事件$A_1,A_2,\ldots$两两互斥，那么
$$P(\bigcup\limits_{i=1}^{+\infty} A_i) = \sum\limits_{i=1}^{+\infty} P(A_i)$$
\end{itemize}
则称$P(A)$为事件$A$的概率。$(\Omega, \mathcal{F}, P)$被称为一个概率空间。
\end{block}

\end{frame}

%------------------------------------------------

\section{随机变量}

%------------------------------------------------

\subsection{随机变量的定义}

%------------------------------------------------

\begin{frame}

\begin{block}{随机变量，Random Variable}
一个随机变量$X$, 指的是一个从{\color{red}样本空间}$\Omega$到实数集$\mathbb{R}$的函数
$$X: \; \Omega \rightarrow \mathbb{R}, \; \omega \mapsto X(\omega)$$
\end{block}

\vspace{1em}
\pause

\begin{remark}
\begin{itemize}
\item 随机变量这个概念将随机试验、随机问题的结果数量化，从而可以用数学的语言加以描述。
\item 随机变量主要分为离散型随机变量与连续型随机变量。{\color{red}注意}！存在既不是离散型也不是连续型的随机变量。
\item 若$g$是一个实变量的实值函数，那么$g(X)$也是一个随机变量。
\end{itemize}
\end{remark}

\end{frame}

%------------------------------------------------

\begin{frame}

\begin{block}{随机变量的例子}
一个赌博游戏，抛2次硬币，都是正面则赢10元，其他输5元，那么样本空间为
$$\Omega = \{\text{正正，反正，正反，反反}\}$$
相应的随机变量就可以定义作
$$X(\omega) = \begin{cases}
10, & \omega = \text{正正} \\
-5, & \text{其他}
\end{cases}$$
\end{block}

\end{frame}

%------------------------------------------------

\subsection{分布函数与密度函数}

%------------------------------------------------

\begin{frame}

\begin{block}{随机变量的（累积）分布函数，(Cumulative) Distribution Function}
设$X$为一个{\color{red}概率空间}$(\Omega, \mathcal{F}, P)$上的随机变量。对任意实数$x\in\mathbb{R}$, 定义
$$F(x) = P(X \le x) = P(\{\omega \ |\ X(\omega) \le x\})$$
为随机变量$X$的（累积）分布函数。此时，称随机变量$X$服从分布$F$, 记作$X \sim F$.
\end{block}

\end{frame}

%------------------------------------------------

\begin{frame}

\begin{block}{分布函数$F(x)$的性质}
\begin{enumerate}
\item 单调性：$F(x)$在$(-\infty, +\infty)$上单调不减
\item 有界性：$0 \le F(x) \le 1$, 并且有
$$F(-\infty) := \lim\limits_{x\to-\infty} F(x) = 0, \quad F(+\infty) := \lim\limits_{x\to+\infty} F(x) = 1$$
\item 右连续性：即对任意$x_0\in\mathbb{R}$, 有
$$F(x_0) = F(x_0+0) := \lim\limits_{x\to x_0+} F(x)$$
\end{enumerate}
\end{block}

\end{frame}

%------------------------------------------------

\begin{frame}

\begin{block}{离散型随机变量的分布 --- 质量函数}
若随机变量$X$的值域是可数个实数$x_1,x_2, \ldots, x_n, \ldots$, 则称其为一个离散型随机变量。特别地，
$$P(X = x_k) = p_k$$
被称作随机变量$X$的分布列，或$X$的概率质量函数（Probability Mass Function, PMF）。于是$X$的分布可以表示为一个表
\begin{center}
\renewcommand{\arraystretch}{1.5}
\begin{tabu} to 0.8\textwidth { X[c] | X[c] | X[c] | X[c] | X[c] | X[c] }
\hline
$X$ & $x_1$ & $x_2$ & $\cdots$ & $x_n$ & $\cdots$ \\
\hline
$p$ & $p_1$ & $p_2$ & $\cdots$ & $p_n$ & $\cdots$ \\
\hline
\end{tabu}
\renewcommand{\arraystretch}{1}
\end{center}
\end{block}

\end{frame}

%------------------------------------------------

\begin{frame}

\begin{block}{连续型随机变量的分布 --- 密度函数}
设随机变量$X$服从分布$F$。若存在非负可积函数$f$，使得对任意$x\in\mathbb{R}$, 有
$$F(x) = \int\limits_{-\infty}^x f(t)dt,$$
则称$X$为连续性随机变量。函数$f$被称为随机变量$X$的概率密度函数（Probability Density Function, PDF）
\end{block}

\end{frame}

%------------------------------------------------

\section{统计量}

%------------------------------------------------

\begin{frame}

\begin{block}{总体与样本}
在统计中，被研究对象的全体被称作总体，可以用随机变量来描述。其组成基本单元被称作个体。$n$个独立，与总体$X$同分布的随机变量$X_1,\ldots,X_n$被称作来自$X$的样本，$n$被称作样本容量。
\end{block}

\vspace{1em}
\pause

\begin{block}{统计量}
设$X_1,\ldots,X_n$是来自总体$X$的样本，$g$是$n$个变元的函数，且不含未知参数，则称$g(X_1,\ldots,X_n)$为一个统计量。常见的统计量有
\begin{itemize}
\item 样本均值：$\overline{X} = \dfrac{X_1 + \cdots + X_n}{n}$
\item 样本方差$S^2 = \dfrac{(X_1-\overline{X})^2 + \cdots + (X_n-\overline{X})^2}{\color{red} n-1}$
\end{itemize}
\end{block}

\end{frame}

% %------------------------------------------------

% \begin{frame}

% \begin{block}{顺序统计量}

% \end{block}

% \end{frame}

%------------------------------------------------

% \begin{frame}

% \begin{block}{Vandermonde行列式的计算}
% \begin{enumerate}
% % \addtocounter{enumi}{0}
% \item 把第$n-1$行乘以$-a_1$加到第$n$行，再把第$n-2$行乘以$-a_1$加到第$n-1$行。按照这种顺序依次向上，直到用第$2$行乘以$-a_1$加到第$1$行。由行列式性质，$V_n$的值不变。
% \[
% V_n = \det \begin{pmatrix}
% 1 & 1 & \cdots & 1 \\
% 0 & a_2-a_1 & \cdots & a_n-a_1 \\
% \vdots & \vdots & & \vdots \\
% 0 & a_2^{n-2}(a_2-a_1) & \cdots & a_n^{n-2}(a_n-a_1) \\
% \end{pmatrix}
% \]
% \end{enumerate}
% \end{block}

% \end{frame}

%------------------------------------------------
% % closing page
% \begin{frame}
% \HUGE{\centerline{\bfseries {\color{gray} The} {\color{zkblue} End}}}

% \pause

% \vspace{1.2em}

% \Large{\centerline{\bfseries {\color{gray}Thank} \quad {\color{zkblue} You}}}

% \end{frame}

%----------------------------------------------------------------------------------------

\end{document}
